\chapter{CONCLUSION AND FUTURE WORKS}
\section{Conclusion}
Our project, GloBus, is a successful venture into building a modern, highly scalable digital marketplace. Moving away from legacy monolithic structures, we adopted the robust MERN stack (MongoDB, Express.js, React.js, Node.js) to deliver a seamless shopping experience for consumers and a powerful management interface for administrators. The primary objective was to create a framework that handles complex e-commerce operations—such as real-time cart state management, AI-driven visual search, and secure payment processing—with absolute efficiency.

At the end of this project lifecycle, it is concluded that we have successfully made efforts on the following points:
\begin{itemize}
    \item Developed a highly responsive frontend using React 19, Vite, and Tailwind CSS v4.
    \item Engineered a secure backend architecture utilizing Express.js and MongoDB.
    \item Integrated advanced AI capabilities, specifically an AI Vision product search module.
    \item Secured user data using Bcrypt password hashing and JWT (JSON Web Tokens).
    \item Implemented the SSLCommerz payment gateway to facilitate local transactions in Bangladesh.
    \item Designed a comprehensive Admin Dashboard for full CRUD operations on products and orders.
    \item Finally, the system was successfully tested across various modules ensuring error-free deployment.
\end{itemize}

\section{Outcomes}
The deployment and testing of GloBus yielded the following positive outcomes:
\begin{itemize}
    \item \textbf{Enhanced Shopping Experience:} Provided an intuitive UI with Dark/Light mode and multi-language support.
    \item \textbf{Secure Transactions:} Successfully bridged the trust gap by integrating verified local payment methods.
    \item \textbf{Technological Advancement:} Showcased how AI (Vision Search and Chatbot) can drastically reduce the effort required for a consumer to find a desired product.
    \item \textbf{Operational Efficiency:} Delivered a turnkey solution for administrators to manage inventory and fulfill orders digitally without requiring third-party software.
\end{itemize}

\section{Current Limitations}
While GloBus is fully operational, there are a few limitations in the current architecture that we aim to overcome in the future:
\begin{itemize}
    \item \textbf{Single-Vendor Architecture:} Currently, GloBus operates on a B2C (Business to Consumer) model where only the central Administrator can upload and manage products. It does not yet support third-party sellers.
    \item \textbf{Manual Logistics Tracking:} While order status can be updated manually in the admin panel, the system is not yet integrated with local courier APIs (e.g., Pathao, RedX) for automated real-time shipping updates.
    \item \textbf{Cash on Delivery (COD) Verification:} Although SSLCommerz handles digital payments flawlessly, enabling a secure COD option in Bangladesh requires strict SMS OTP verification to prevent fake orders, which is currently missing.
\end{itemize}

\section{Future Works}
GloBus has laid a powerful foundation. If the platform scales, several vital features can be integrated to compete directly with massive e-commerce giants.
\begin{itemize}
    \item \textbf{Multi-Vendor Marketplace Transition:} In the future, we aim to transition the database schema and routing to allow independent sellers to open their own "shops" within GloBus, complete with a vendor-specific dashboard and commission tracking.
    \item \textbf{Mobile Application Development:} Utilizing our existing Node.js backend, we plan to develop a cross-platform mobile application using React Native to capture the massive mobile-first consumer base.
    \item \textbf{AI Recommendation Engine:} We intend to implement a Machine Learning model that tracks user behavior (clicks, cart additions, previous purchases) to provide highly personalized product recommendations on the home page.
    \item \textbf{Logistics \& Delivery API Integration:} Integrate third-party courier services directly into the backend so that when an order is marked as "Shipped", the tracking ID automatically fetches live geolocation data for the user.
    \item \textbf{Advanced Security:} Implement Two-Factor Authentication (2FA) via SMS or Authenticator apps for both standard users and administrators to elevate platform security to enterprise standards.
\end{itemize}
